\documentclass[a4paper, 12pt]{article}

\usepackage[brazilian]{babel}
\usepackage[utf8]{inputenc}
\usepackage{amsmath, amssymb}
\usepackage{indentfirst}
\usepackage{graphicx}
\usepackage{multicol}
\usepackage{url}
\usepackage{hyperref}
\usepackage{booktabs}
\usepackage{longtable}
\usepackage{tabularx}
\usepackage{enumitem}
\usepackage{float}
\usepackage{xcolor}
\usepackage{tcolorbox}
\usepackage{setspace}
\usepackage{fancyhdr}
\usepackage{titlesec}
\usepackage[top=2.0cm, bottom=2.0cm, left=2.5cm, right=2.0cm]{geometry}

% ── Cores do projeto ────────────────────────────────────────────────────────
\definecolor{beegeAmbar}{HTML}{F59E0B}
\definecolor{beegePurple}{HTML}{7C3AED}
\definecolor{beegeGray}{HTML}{374151}

% ── Caixa de destaque ──────────────────────────────────────────────────────
\newtcolorbox{destaque}[1][]{
  colback=beegeAmbar!8,
  colframe=beegeAmbar!80!black,
  fonttitle=\bfseries,
  title={#1},
  arc=3pt,
  boxrule=0.8pt
}

\hypersetup{
    colorlinks=true,
    linkcolor=beegePurple,
    citecolor=beegePurple,
    urlcolor=beegePurple
}

\setstretch{1.5}

% ── Cabecalho e rodape ─────────────────────────────────────────────────────
\pagestyle{fancy}
\fancyhf{}
\fancyhead[L]{\small\textcolor{beegeGray}{BeeGe --- Salto Tecnol\'{o}gico}}
\fancyhead[R]{\small\textcolor{beegeGray}{NaPr\'{a}tica 2026}}
\fancyfoot[C]{\thepage}
\renewcommand{\headrulewidth}{0.4pt}

\begin{document}
\hypersetup{pageanchor=false}

% % % % % % % % % % % % % % % CAPA % % % % % % % % % % % % % % %
\begin{titlepage}
	\begin{center}
		\textbf{\large UNIVERSIDADE DE BRAS\'{I}LIA -- UnB}\\
		\textbf{\large FACULDADE DE CI\^{E}NCIAS E TECNOLOGIAS EM ENGENHARIA -- FCTE}\\
		\textbf{\large CURSO DE ENGENHARIA DE SOFTWARE}\\
		\large PROGRAMA NaPr\'{a}tica --- M\'{O}DULO DE EMPREENDEDORISMO\\
		\vspace{35pt}
		
		\textbf{\large GABRIEL MONTEIRO}\\
		\vspace{6pt}
		{\normalsize
		Engenheiro de Software em Forma\c{c}\~{a}o $\cdot$ Tutor de C\'{a}lculo~1 na UnB
		}
		
		\vspace{70pt}
		
		\textbf{\LARGE{SALTO TECNOL\'{O}GICO: BeeGe}}\\
		\vspace{12pt}
		\textbf{\Large{Plataforma Mobile Gamificada de Revis\~{a}o de C\'{a}lculo~1 com Arquitetura SaaS Multi-Institucional}}\\
		
		\vspace{\fill}
		
		Gama -- DF\\
		Outubro de 2026
	\end{center}
\end{titlepage}

% % % % % % % % % % % % % FOLHA DE ROSTO % % % % % % % % % % % %
\begin{titlepage}
	\begin{center}
		\textbf{\large UNIVERSIDADE DE BRAS\'{I}LIA -- UnB}\\
		\textbf{\large FACULDADE DE CI\^{E}NCIAS E TECNOLOGIAS EM ENGENHARIA -- FCTE}\\
		\textbf{\large ENGENHARIA DE SOFTWARE}\\
		\vspace{35pt}
		
		\textbf{\large GABRIEL MONTEIRO}\\
		\vspace{6pt}
		{\normalsize
		23 anos $\cdot$ Tutor de C\'{a}lculo~1 h\'{a} mais de 3 anos
		}
		
		\vspace{50pt}
		
		\textbf{\LARGE{SALTO TECNOL\'{O}GICO: BeeGe}}\\
		\vspace{12pt}
		\textbf{\Large{Plataforma Mobile Gamificada de Revis\~{a}o de C\'{a}lculo~1 com Arquitetura SaaS Multi-Institucional}}\\
	\end{center}
	
	\vspace{30pt}
	
	\begin{flushright}
		\begin{minipage}{0.58\textwidth}
			\small
			Documento de apresenta\c{c}\~{a}o do \textbf{Salto Tecnol\'{o}gico} desenvolvido como entreg\'{a}vel do programa \textbf{NaPr\'{a}tica --- M\'{o}dulo de Empreendedorismo}. O projeto consiste no desenvolvimento de um MVP (\textit{Minimum Viable Product}) de uma plataforma mobile e web para revis\~{a}o gamificada de C\'{a}lculo~1, concebido em 30 dias com potencial de transforma\c{c}\~{a}o em produto SaaS.
			\vspace{10pt}
			
			\textbf{Autor:} Gabriel Monteiro\\
			\textbf{Produto:} BeeGe --- Plataforma de Revis\~{a}o de C\'{a}lculo~1\\
			\textbf{Per\'{i}odo:} 20 de Setembro a 19 de Outubro de 2026\\
			\textbf{Reposit\'{o}rio:} \url{https://github.com/gabriel-monteiro/beege}
		\end{minipage}
	\end{flushright}
	
	\vspace{\fill}
	
	\begin{center}
		Gama -- DF\\
		2026
	\end{center}
\end{titlepage}

\tableofcontents
\thispagestyle{empty}

\newpage
\hypersetup{pageanchor=true}
\pagenumbering{arabic}

% % % % % % % % % % % % SE\c{C}\~{O}ES DO ARTIGO % % % % % % % % % % %
\section{Resumo}
Este trabalho apresenta o \textbf{BeeGe}, uma plataforma mobile e web voltada para a revis\~{a}o gamificada de C\'{a}lculo~1 no ensino superior, desenvolvida como Salto Tecnol\'{o}gico no programa NaPr\'{a}tica. O projeto foi concebido, projetado e implementado em 30 dias por Gabriel Monteiro, estudante de Engenharia de Software e tutor de C\'{a}lculo~1 na Universidade de Bras\'{i}lia (UnB). O BeeGe organiza conte\'{u}dos por universidade (modelo SaaS multi-institucional), integra renderiza\c{c}\~{a}o de f\'{o}rmulas matem\'{a}ticas via KaTeX, oferece exerc\'{i}cios com resolu\c{c}\~{a}o anal\'{i}tica passo a passo e conta com um f\'{o}rum acad\^{e}mico moderado. O MVP entregue contempla 15 requisitos funcionais, 8 n\~{a}o-funcionais, 18 hist\'{o}rias de usu\'{a}rio implementadas, 24 quest\~{o}es com gabarito comentado e 12 documentos de engenharia de software. A arquitetura adota Ionic Framework com React~19, Supabase (PostgreSQL com RLS) e Capacitor para compila\c{c}\~{a}o nativa, operando em Android, iOS e web a partir de base \'{u}nica de c\'{o}digo.

\vspace{8pt}
\noindent\textbf{Palavras-chave:} Empreendedorismo, EdTech, C\'{a}lculo~1, SaaS, MVP, Ionic, Supabase, KaTeX, Gamifica\c{c}\~{a}o.

% ================================================================
\section{Introdu\c{c}\~{a}o}
\label{sec:introducao}

O ensino e a aprendizagem de C\'{a}lculo~1 no ensino superior historicamente registram \'{i}ndices elevados de reten\c{c}\~{a}o e reprova\c{c}\~{a}o em universidades de todo o Brasil. Estudantes frequentemente relatam que livros did\'{a}ticos s\~{a}o excessivamente extensos --- muitos com mais de 800 p\'{a}ginas \cite{stewart2016} --- e que videoaulas dispon\'{i}veis na internet n\~{a}o seguem o estilo de resolu\c{c}\~{a}o ou a profundidade exigida nas provas de sua institui\c{c}\~{a}o espec\'{i}fica.

Diante dessa realidade, nasceu o \textbf{BeeGe}: uma plataforma mobile e web focada na revis\~{a}o e fixa\c{c}\~{a}o de C\'{a}lculo~1, estruturada como um servi\c{c}o multi-institucional (SaaS). O projeto foi desenvolvido como \textbf{Salto Tecnol\'{o}gico} no programa NaPr\'{a}tica, unindo tr\^{e}s verte\'{n}tes: a experi\^{e}ncia do autor como tutor de C\'{a}lculo~1 na Universidade de Bras\'{i}lia (UnB) h\'{a} mais de tr\^{e}s anos, a forma\c{c}\~{a}o em Engenharia de Software e o desafio pessoal de empreender.

O resultado \'{e} um \textit{Minimum Viable Product} (MVP) \cite{ries2011} funcional, desenvolvido em 30 dias --- de 20 de setembro a 19 de outubro de 2026 --- com potencial de se transformar em um produto SaaS capaz de atender universidades e faculdades em todo o pa\'{i}s.

\textbf{Estrutura do artigo.} A Se\c{c}\~{a}o~\ref{sec:motivacao} apresenta a motiva\c{c}\~{a}o pessoal e o contexto; a Se\c{c}\~{a}o~\ref{sec:problema} detalha o problema identificado; a Se\c{c}\~{a}o~\ref{sec:solucao} descreve a solu\c{c}\~{a}o proposta; a Se\c{c}\~{a}o~\ref{sec:arquitetura} apresenta a arquitetura t\'{e}cnica; a Se\c{c}\~{a}o~\ref{sec:funcionalidades} enumera as funcionalidades implementadas; a Se\c{c}\~{a}o~\ref{sec:metodologia} descreve a metodologia adotada; a Se\c{c}\~{a}o~\ref{sec:cronograma} detalha o cronograma de desenvolvimento; a Se\c{c}\~{a}o~\ref{sec:qualidade} aborda qualidade e testes; a Se\c{c}\~{a}o~\ref{sec:backlog} apresenta o backlog; a Se\c{c}\~{a}o~\ref{sec:negocio} discute o modelo de neg\'{o}cio; a Se\c{c}\~{a}o~\ref{sec:resultados} consolida os resultados; a Se\c{c}\~{a}o~\ref{sec:documentacao} lista a documenta\c{c}\~{a}o produzida; e a Se\c{c}\~{a}o~\ref{sec:conclusao} re\'{u}ne as conclus\~{o}es.

% ================================================================
\section{Motiva\c{c}\~{a}o Pessoal e Contexto}
\label{sec:motivacao}

\begin{destaque}[Depoimento do Autor]
Me chamo Gabriel Monteiro, tenho 23 anos e decidi fazer esse projeto para criar a coragem de empreender. No m\'{o}dulo~01, participei animado das palestras e ensinamentos e tive a ideia de me desafiar a fazer algo que eu sempre quis, mas n\~{a}o tinha coragem: \textbf{empreender}.

Quando eu estava l\'{a}, ouvi outros alunos falando de como estavam sofrendo com a mat\'{e}ria de C\'{a}lculo~1 na faculdade. At\'{e} ent\~{a}o, isso era normal para mim, dado que dou tutoria de C\'{a}lculo~1 na UnB h\'{a} mais de 3 anos, mas eu nunca havia pensado nisso como uma oportunidade de neg\'{o}cio.

Unindo esse desafio que nasceu, ser um engenheiro de software em forma\c{c}\~{a}o e t\~{a}o ligado \`{a} \'{a}rea de educa\c{c}\~{a}o e matem\'{a}tica, me desafiei a fazer esse projeto. Quando tivesse um MVP consolidado, eu poderia abrir a minha empresa e transformar esse projeto em um SaaS para que diversas universidades e faculdades possam modelar as quest\~{o}es da maneira que acharem mais adequadas, j\'{a} que isso muda de grade para grade curricular.

Dessa maneira nasceu o \textbf{BeeGe}, uma plataforma mobile capaz de ensinar C\'{a}lculo~1 de forma gamificada, mas com rigor te\'{o}rico. Durante o desenvolvimento, falei com professores de C\'{a}lculo~1 para entender o que era mais fundamental e fazer o meu brainstorm.

A ideia, por ter apenas 30 dias, era criar um MVP capaz de fazer as funcionalidades b\'{a}sicas: ter a parte de \textit{admin} para a cria\c{c}\~{a}o de trilhas e quest\~{o}es, ter a parte de estudantes para acessar as trilhas, conte\'{u}dos e quest\~{o}es, al\'{e}m de ter uma ofensiva para que os alunos se mantenham engajados no aprendizado.

Dessa feita, acredito que impactaria de maneira muito positiva todos os alunos de C\'{a}lculo~1, ou que quisessem revisar a base matem\'{a}tica, e ainda estaria entrando no mundo de empreendedorismo que sempre tive medo de arriscar.
\end{destaque}

% ================================================================
\section{O Problema}
\label{sec:problema}

O ensino de C\'{a}lculo~1 apresenta desafios estruturais que afetam estudantes em todo o pa\'{i}s. As principais dificuldades relatadas s\~{a}o:

\begin{itemize}[leftmargin=2cm]
  \item Livros did\'{a}ticos excessivamente extensos e abstratos, com abordagens generalistas que n\~{a}o refletem o estilo de cobran\c{c}a de cada institui\c{c}\~{a}o.
  \item Videoaulas dispon\'{i}veis na web que n\~{a}o seguem a profundidade exigida nas provas da universidade espec\'{i}fica do estudante.
  \item Grupos de mensagens que acumulam conversas paralelas, dificultando a localiza\c{c}\~{a}o de d\'{u}vidas j\'{a} respondidas.
  \item Materiais de monitoria distribu\'{i}dos de forma fragmentada em PDFs que se perdem nos dispositivos.
\end{itemize}

Existe, portanto, a demanda por uma solu\c{c}\~{a}o orientada \`{a} revis\~{a}o que combine: (1)~segmenta\c{c}\~{a}o por universidade, respeitando a ementa local; (2)~renderiza\c{c}\~{a}o digital de nota\c{c}\~{a}o matem\'{a}tica em alta qualidade; (3)~exerc\'{i}cios com passo a passo anal\'{i}tico dos gabaritos; e (4)~ambiente colaborativo de d\'{u}vidas com modera\c{c}\~{a}o ativa.

% ================================================================
\section{A Solu\c{c}\~{a}o: BeeGe}
\label{sec:solucao}

\subsection{Vis\~{a}o do Produto}

O \textbf{BeeGe} \'{e} um aplicativo mobile e web focado na revis\~{a}o e fixa\c{c}\~{a}o de C\'{a}lculo~1, estruturado como um servi\c{c}o multi-institucional (SaaS). Diferentemente de cursos em v\'{i}deo gen\'{e}ricos ou plataformas abertas de quest\~{o}es, o BeeGe organiza o conte\'{u}do a partir da ementa espec\'{i}fica da institui\c{c}\~{a}o do estudante \cite{ionic2024}.

\begin{table}[H]
\centering
\caption{Posicionamento estrat\'{e}gico do produto}
\label{tab:posicionamento}
\vspace{4pt}
\small
\begin{tabularx}{\textwidth}{l X}
\toprule
\textbf{Dimens\~{a}o} & \textbf{Especifica\c{c}\~{a}o} \\
\midrule
\textbf{Para} & Estudantes de gradua\c{c}\~{a}o matriculados na disciplina de C\'{a}lculo~1. \\
\textbf{Que buscam} & Material de revis\~{a}o objetivo e focado no modelo de prova de sua universidade. \\
\textbf{O BeeGe \'{e}} & Um aplicativo mobile (Android/iOS) e plataforma web. \\
\textbf{Que oferece} & Teoria sintetizada, exerc\'{i}cios corrigidos instantaneamente e canal de d\'{u}vidas. \\
\textbf{Diferente de} & Cursos longos em v\'{i}deo e manuais tradicionais extensos. \\
\textbf{Diferencial} & Adapta\c{c}\~{a}o por universidade, KaTeX nativo e suporte da mascote Limia. \\
\bottomrule
\end{tabularx}
\end{table}

\subsection{Identidade Visual e Mascote (Limia)}

A mascote oficial do projeto \'{e} a abelha \textbf{Limia}:

\begin{itemize}[leftmargin=2cm]
  \item \textbf{Caracter\'{i}sticas Visuais:} Abelha estilizada com \'{o}culos redondos pretos e gravata borboleta roxa contendo s\'{i}mbolos matem\'{a}ticos gravados ($\int$, $\lim$, $\partial$).
  \item \textbf{Simbolismo:} Representa o rigor acad\^{e}mico aliado \`{a} dilig\^{e}ncia e ao trabalho em colmeia (aprendizagem colaborativa).
  \item \textbf{Papel no Sistema:} Atua como elemento de apoio pedag\'{o}gico, apresentando dicas conceituais, observa\c{c}\~{o}es sobre teoremas e mensagens de orienta\c{c}\~{a}o.
\end{itemize}

\subsection{P\'{u}blico-Alvo (Personas)}

O sistema atende tr\^{e}s perfis principais de usu\'{a}rios, levantados durante a fase de pesquisa:

\begin{table}[H]
\centering
\caption{Perfis de usu\'{a}rios do BeeGe}
\label{tab:personas}
\vspace{4pt}
\small
\begin{tabularx}{\textwidth}{l l l l X}
\toprule
\textbf{Persona} & \textbf{Idade} & \textbf{Papel} & \textbf{Dispositivo} & \textbf{Necessidade Principal} \\
\midrule
Lucas (Calouro) & 18 & Aluno & Smartphone & Compreender limites e derivadas no ritmo da UnB \\
Mariana (Veterana) & 21 & Aluno & Smartphone & Revis\~{a}o r\'{a}pida e pr\'{a}tica para recupera\c{c}\~{a}o \\
Rafael (Monitor) & 23 & Admin & Desktop & Centralizar exerc\'{i}cios e moderar d\'{u}vidas \\
\bottomrule
\end{tabularx}
\end{table}

% ================================================================
\section{Arquitetura e Tecnologias}
\label{sec:arquitetura}

O sistema BeeGe adota o padr\~{a}o arquitetural \textbf{Cliente-Servidor baseado em BaaS} (\textit{Backend as a Service}) \cite{supabase2024}. A interface opera como uma aplica\c{c}\~{a}o de p\'{a}gina \'{u}nica (SPA) multiplataforma que se comunica de maneira ass\'{i}ncrona com a infraestrutura gerenciada pelo Supabase.

\subsection{Stack Tecnol\'{o}gico}

\begin{table}[H]
\centering
\caption{Tecnologias adotadas no MVP}
\label{tab:tecnologias}
\vspace{4pt}
\small
\begin{tabularx}{\textwidth}{l l X l}
\toprule
\textbf{Camada} & \textbf{Tecnologia} & \textbf{Justificativa} & \textbf{Custo} \\
\midrule
Frontend & Ionic + React~19 & Interface responsiva unificada para web e mobile \cite{ionic2024} & Gratuito \\
Mobile & Capacitor & Builds Android e iOS sem duplica\c{c}\~{a}o de base \cite{capacitor2024} & Gratuito \\
Backend & Supabase (PostgreSQL) & Servi\c{c}o gerenciado com RLS, Auth e fun\c{c}\~{o}es \cite{supabase2024} & Gratuito \\
F\'{o}rmulas & KaTeX & Renderiza\c{c}\~{a}o performatica no cliente \cite{katex2024} & Gratuito \\
Documenta\c{c}\~{a}o & MkDocs Material & Gerador est\'{a}tico versionado com o c\'{o}digo & Gratuito \\
Testes & Vitest + Cypress & Testes unit\'{a}rios, de componentes e E2E & Gratuito \\
Linguagem & TypeScript~5.9 & Tipagem estrita em tempo de compila\c{c}\~{a}o & Gratuito \\
\bottomrule
\end{tabularx}
\end{table}

\subsection{Camada Cliente (\texttt{beege-app})}

\begin{itemize}[leftmargin=2cm]
  \item \textbf{Framework de Interface:} React~19 com a biblioteca de componentes gr\'{a}ficos do Ionic Framework.
  \item \textbf{Navega\c{c}\~{a}o:} \texttt{@ionic/react-router} integrado ao React Router~6.
  \item \textbf{Renderiza\c{c}\~{a}o Matem\'{a}tica:} Componente dedicado \texttt{MathRenderer} que encapsula o motor KaTeX.
  \item \textbf{Gerenciamento de Estado:} Context API (\texttt{AuthContext}) para sess\~{a}o e universidade ativa.
  \item \textbf{Empacotamento Mobile:} Capacitor para compila\c{c}\~{a}o em Android (APK) e iOS (IPA).
\end{itemize}

\subsection{Camada Backend (\texttt{beege-api})}

\begin{itemize}[leftmargin=2cm]
  \item \textbf{Autentica\c{c}\~{a}o:} Supabase Auth com tokens JWT e controle de sess\~{a}o.
  \item \textbf{Persist\^{e}ncia:} PostgreSQL~15 com extens\~{o}es como \texttt{uuid-ossp}.
  \item \textbf{Seguran\c{c}a (RLS):} Regras granulares por linha para isolamento de dados.
  \item \textbf{Migra\c{c}\~{o}es:} Arquivos SQL versionados com \textit{timestamps}.
\end{itemize}

\subsection{Padr\~{o}es de Projeto Utilizados}

\begin{itemize}[leftmargin=2cm]
  \item \textbf{Component-Based Architecture:} Modulariza\c{c}\~{a}o em componentes funcionais reutiliz\'{a}veis com React Hooks \cite{pressman2021}.
  \item \textbf{Service Layer Pattern:} Centraliza\c{c}\~{a}o de requisi\c{c}\~{o}es na camada \texttt{DataService}.
  \item \textbf{Observer Pattern:} \texttt{onAuthStateChange} do Supabase Auth para reatividade do estado de login.
\end{itemize}

% ================================================================
\section{Funcionalidades do MVP}
\label{sec:funcionalidades}

\subsection{Escopo Funcional Implementado}

\begin{longtable}{l p{8cm} l}
\caption{M\'{o}dulos e funcionalidades implementadas no MVP} \\
\toprule
\textbf{M\'{o}dulo} & \textbf{Funcionalidade} & \textbf{Status} \\
\midrule
\endfirsthead
\toprule
\textbf{M\'{o}dulo} & \textbf{Funcionalidade} & \textbf{Status} \\
\midrule
\endhead
Autentica\c{c}\~{a}o & Cadastro com email e senha & Implementado \\
Autentica\c{c}\~{a}o & Login persistente e recupera\c{c}\~{a}o de credenciais & Implementado \\
Multi-Universidade & Sele\c{c}\~{a}o da institui\c{c}\~{a}o e filtragem autom\'{a}tica & Implementado \\
Conte\'{u}do & Listagem sequencial dos m\'{o}dulos de C\'{a}lculo~1 & Implementado \\
Renderiza\c{c}\~{a}o & Suporte nativo a sintaxe \LaTeX{} com KaTeX & Implementado \\
Quest\~{o}es & Exerc\'{i}cios de m\'{u}ltipla escolha com gabarito comentado & Implementado \\
Quest\~{o}es & Avalia\c{c}\~{a}o imediata com explica\c{c}\~{a}o anal\'{i}tica & Implementado \\
Progresso & Hist\'{o}rico persistente via \texttt{ProgressService} & Implementado \\
Progresso & Streak, taxa de acertos, meta di\'{a}ria e conquistas & Implementado \\
F\'{o}rum & Cria\c{c}\~{a}o de postagens com suporte a \LaTeX{} & Implementado \\
F\'{o}rum & Respostas encadeadas e modera\c{c}\~{a}o pr\'{e}via & Implementado \\
Admin & Cadastro e edi\c{c}\~{a}o de t\'{o}picos te\'{o}ricos & Implementado \\
Admin & Cadastro de quest\~{o}es com alternativas e gabarito & Implementado \\
Admin & Fila de modera\c{c}\~{a}o de postagens do f\'{o}rum & Implementado \\
Mascote & Dicas contextuais da Limia em todas as telas & Implementado \\
Seguran\c{c}a & Rotas protegidas e guarda de acesso administrativo & Implementado \\
\bottomrule
\end{longtable}

\subsection{Requisitos Funcionais (RF)}

\begin{longtable}{l p{9cm} l}
\caption{Requisitos funcionais especificados e atendidos} \\
\toprule
\textbf{ID} & \textbf{Descri\c{c}\~{a}o} & \textbf{Prioridade} \\
\midrule
\endfirsthead
\toprule
\textbf{ID} & \textbf{Descri\c{c}\~{a}o} & \textbf{Prioridade} \\
\midrule
\endhead
RF01 & Cadastro de novos usu\'{a}rios com nome, email e senha & Essencial \\
RF02 & Autentica\c{c}\~{a}o e gerenciamento de sess\~{o}es & Essencial \\
RF03 & Sele\c{c}\~{a}o e altern\^{a}ncia da universidade ativa & Essencial \\
RF04 & Listagem sequencial dos m\'{o}dulos te\'{o}ricos & Essencial \\
RF05 & Exibi\c{c}\~{a}o de conte\'{u}do te\'{o}rico com f\'{o}rmulas em KaTeX & Essencial \\
RF06 & Exerc\'{i}cios de m\'{u}ltipla escolha por t\'{o}pico & Essencial \\
RF07 & Avalia\c{c}\~{a}o imediata (acerto/erro) & Essencial \\
RF08 & Explica\c{c}\~{a}o anal\'{i}tica com nota\c{c}\~{a}o \LaTeX{} ap\'{o}s resposta & Essencial \\
RF09 & Registro de pontua\c{c}\~{a}o e progresso & Essencial \\
RF10 & Postagens no f\'{o}rum com suporte a \LaTeX{} & Importante \\
RF11 & Respostas e coment\'{a}rios no f\'{o}rum & Importante \\
RF12 & Cadastro e edi\c{c}\~{a}o de m\'{o}dulos te\'{o}ricos (Admin) & Essencial \\
RF13 & Cadastro de quest\~{o}es com gabarito (Admin) & Essencial \\
RF14 & Modera\c{c}\~{a}o de d\'{u}vidas do f\'{o}rum (Admin) & Essencial \\
RF15 & Orienta\c{c}\~{o}es conceituais da mascote Limia & Importante \\
\bottomrule
\end{longtable}

\subsection{Requisitos N\~{a}o-Funcionais (RNF)}

\begin{table}[H]
\centering
\caption{Requisitos n\~{a}o-funcionais especificados}
\label{tab:rnf}
\vspace{4pt}
\footnotesize
\begin{tabular}{l p{6cm} l l}
\toprule
\textbf{ID} & \textbf{Descri\c{c}\~{a}o} & \textbf{Categoria} & \textbf{Prioridade} \\
\midrule
RNF01 & Opera\c{c}\~{a}o em Android e iOS via Capacitor/Ionic & Portabilidade & Essencial \\
RNF02 & Idioma padr\~{a}o: Portugu\^{e}s do Brasil (pt-BR) & Localiza\c{c}\~{a}o & Essencial \\
RNF03 & Renderiza\c{c}\~{a}o fluida de f\'{o}rmulas com KaTeX & Usabilidade & Essencial \\
RNF04 & PostgreSQL no Supabase com RLS & Arquitetura & Essencial \\
RNF05 & Criptografia de senhas e tokens & Seguran\c{c}a & Essencial \\
RNF06 & Camada de dados local para execu\c{c}\~{a}o offline & Confiabilidade & Essencial \\
RNF07 & Interface responsiva de 360px a desktops & Usabilidade & Essencial \\
RNF08 & Carregamento inicial abaixo de 1 segundo & Desempenho & Importante \\
\bottomrule
\end{tabular}
\end{table}

% ================================================================
\section{Metodologia de Desenvolvimento}
\label{sec:metodologia}

\subsection{Scrum Adaptado}

Adotou-se uma varia\c{c}\~{a}o pragm\'{a}tica do framework Scrum \cite{schwaber2020} com ciclos semanais, adequada ao prazo de 30 dias:

\begin{itemize}[leftmargin=2cm]
  \item \textbf{Sprints de 1 semana:} Itera\c{c}\~{o}es curtas com foco em entregas incrementais.
  \item \textbf{Planejamento de Sprint:} Sele\c{c}\~{a}o de hist\'{o}rias do backlog priorit\'{a}rias.
  \item \textbf{Revis\~{a}o e Retrospectiva:} Avalia\c{c}\~{a}o dos entreg\'{a}veis ao t\'{e}rmino de cada sprint.
  \item \textbf{Kanban Integrado:} Gest\~{a}o visual: \texttt{Backlog} $\rightarrow$ \texttt{To~Do} $\rightarrow$ \texttt{In~Progress} $\rightarrow$ \texttt{Review} $\rightarrow$ \texttt{Done}.
\end{itemize}

\subsection{Pr\'{a}ticas de Engenharia de Software}

\begin{itemize}[leftmargin=2cm]
  \item \textbf{Arquitetura Modular:} Separa\c{c}\~{a}o entre \texttt{beege-docs}, \texttt{beege-app} e \texttt{beege-api} \cite{pressman2021}.
  \item \textbf{Conventional Commits:} Padroniza\c{c}\~{a}o do hist\'{o}rico no Git.
  \item \textbf{Resili\^{e}ncia Local:} Camada de dados \textit{mockada} para opera\c{c}\~{a}o offline.
  \item \textbf{Design-First com KaTeX:} Fidelidade visual de equa\c{c}\~{o}es desde a concep\c{c}\~{a}o.
\end{itemize}

% ================================================================
\section{Cronograma de Desenvolvimento}
\label{sec:cronograma}

O MVP foi desenvolvido em 5 sprints semanais, totalizando 30 dias:

\begin{table}[H]
\centering
\caption{Cronograma de sprints executadas}
\label{tab:cronograma}
\vspace{4pt}
\small
\begin{tabularx}{\textwidth}{l l X}
\toprule
\textbf{Per\'{i}odo} & \textbf{Marco} & \textbf{Principais Entreg\'{a}veis} \\
\midrule
20/09 -- 23/09 & Sprint~0 & Defini\c{c}\~{a}o de arquitetura, reposit\'{o}rios e documenta\c{c}\~{a}o base \\
24/09 -- 30/09 & Sprint~1 & Modelagem de dados, Supabase e fluxo de autentica\c{c}\~{a}o \\
01/10 -- 07/10 & Sprint~2 & Painel administrativo, t\'{o}picos e resolu\c{c}\~{a}o de quest\~{o}es \\
08/10 -- 13/10 & Sprint~3 & F\'{o}rum, modera\c{c}\~{a}o, integra\c{c}\~{a}o KaTeX e refinamento de UI \\
14/10 -- 19/10 & Sprint~4 & Valida\c{c}\~{a}o, testes, builds e entrega final do MVP \\
\bottomrule
\end{tabularx}
\end{table}

\subsection{Destaques da Sprint Final (Sprint~4)}

\begin{itemize}[leftmargin=2cm]
  \item Expans\~{a}o para 24 quest\~{o}es com gabarito e explica\c{c}\~{a}o anal\'{i}tica em KaTeX (meta de 20+ superada).
  \item Persist\^{e}ncia de progresso via \texttt{ProgressService} (streak, taxa de acertos, meta di\'{a}ria e conquistas).
  \item Dicas contextuais da mascote Limia em Dashboard, Teoria, Pr\'{a}tica, F\'{o}rum e Perfil.
  \item Seguran\c{c}a de rotas com \texttt{ProtectedRoute} e guarda de acesso administrativo.
  \item 17 testes unit\'{a}rios automatizados com 100\% de sucesso.
  \item Build de produ\c{c}\~{a}o bem-sucedido e sincroniza\c{c}\~{a}o com Capacitor.
\end{itemize}

% ================================================================
\section{Qualidade de Software e Testes}
\label{sec:qualidade}

\subsection{Ferramental}

\begin{itemize}[leftmargin=2cm]
  \item \textbf{TypeScript:} Verifica\c{c}\~{a}o estrita de tipos (\texttt{tsc --noEmit}).
  \item \textbf{ESLint:} Padroniza\c{c}\~{a}o e preven\c{c}\~{a}o de antipadr\~{o}es React/Ionic.
  \item \textbf{Vitest + React Testing Library:} Testes unit\'{a}rios e de componentes \cite{basili1994gqm}.
  \item \textbf{Cypress:} Testes E2E simulando navega\c{c}\~{a}o completa do usu\'{a}rio.
  \item \textbf{CI/CD:} Valida\c{c}\~{a}o autom\'{a}tica via GitHub Actions a cada Pull Request.
\end{itemize}

\subsection{M\'{e}tricas Alvo}

\begin{table}[H]
\centering
\caption{M\'{e}tricas de qualidade do MVP}
\label{tab:metricas}
\vspace{4pt}
\small
\begin{tabularx}{\textwidth}{l l X}
\toprule
\textbf{M\'{e}trica} & \textbf{Alvo} & \textbf{Mecanismo de Aferi\c{c}\~{a}o} \\
\midrule
Sucesso de Compila\c{c}\~{a}o & 100\% & Pipeline CI / \texttt{npm run build} \\
Tempo de Carregamento & $< 1{,}5\text{s}$ em 4G & Google Lighthouse / DevTools \\
Cobertura de Regras de Neg\'{o}cio & $\geq 80\%$ & Vitest Coverage Report \\
Erros de Renderiza\c{c}\~{a}o KaTeX & Zero & Testes de regress\~{a}o visual \\
\bottomrule
\end{tabularx}
\end{table}

% ================================================================
\section{Backlog: Hist\'{o}rias de Usu\'{a}rio}
\label{sec:backlog}

As 18 hist\'{o}rias de usu\'{a}rio foram organizadas em 6 \'{e}picos funcionais no formato \textit{``Como [ator], desejo [funcionalidade] para que [benef\'{i}cio]''}:

\begin{enumerate}[leftmargin=2cm]
  \item \textbf{\'{E}pico 1 --- Autentica\c{c}\~{a}o:} Cadastro, login e encerramento de sess\~{a}o (US01--US03).
  \item \textbf{\'{E}pico 2 --- Multi-Universidade:} Sele\c{c}\~{a}o e altern\^{a}ncia de universidade ativa (US04--US05).
  \item \textbf{\'{E}pico 3 --- Conte\'{u}do Te\'{o}rico:} M\'{o}dulos, teoria com KaTeX e dicas da Limia (US06--US08).
  \item \textbf{\'{E}pico 4 --- Pr\'{a}tica:} Resolu\c{c}\~{a}o, feedback imediato e progresso (US09--US12).
  \item \textbf{\'{E}pico 5 --- F\'{o}rum Acad\^{e}mico:} D\'{u}vidas, respostas e visualiza\c{c}\~{a}o moderada (US13--US15).
  \item \textbf{\'{E}pico 6 --- Painel Administrativo:} T\'{o}picos, quest\~{o}es e modera\c{c}\~{a}o (US16--US18).
\end{enumerate}

\textbf{Todas as 18 hist\'{o}rias foram implementadas e entregues no MVP.}

% ================================================================
\section{Modelo de Neg\'{o}cio e Vis\~{a}o de Futuro}
\label{sec:negocio}

\subsection{De MVP a SaaS}

O BeeGe foi projetado desde a concep\c{c}\~{a}o como uma plataforma \textbf{multi-inquilino} (multi-tenant) \cite{sommerville2019}. Cada universidade opera como um \textit{tenant} isolado, com seus pr\'{o}prios t\'{o}picos, quest\~{o}es e f\'{o}rum. Essa arquitetura viabiliza a transforma\c{c}\~{a}o do MVP em produto SaaS:

\begin{itemize}[leftmargin=2cm]
  \item Universidades modelam quest\~{o}es e trilhas de acordo com suas grades curriculares.
  \item Monitores e professores criam e moderam conte\'{u}do com autonomia.
  \item Escala horizontal --- novas institui\c{c}\~{o}es s\~{a}o adicionadas sem altera\c{c}\~{a}o no c\'{o}digo.
\end{itemize}

\subsection{Extensibilidade}

A arquitetura permite receber disciplinas subsequentes: C\'{a}lculo~2 (integrais m\'{u}ltiplas), C\'{a}lculo~3 (c\'{a}lculo vetorial), \'{A}lgebra Linear e Equa\c{c}\~{o}es Diferenciais.

\subsection{Potencial de Impacto}

\begin{destaque}[Impacto Esperado]
O BeeGe tem potencial para impactar positivamente milhares de estudantes de C\'{a}lculo~1 em todo o Brasil, oferecendo uma ferramenta de revis\~{a}o acess\'{i}vel, gamificada e alinhada com a realidade de cada institui\c{c}\~{a}o. Al\'{e}m disso, representa a concretiza\c{c}\~{a}o de um sonho empreendedor que nasceu dentro do programa NaPr\'{a}tica.
\end{destaque}

% ================================================================
\section{Resultados Alcan\c{c}ados}
\label{sec:resultados}

\begin{table}[H]
\centering
\caption{M\'{e}tricas de sucesso do MVP}
\label{tab:resultados}
\vspace{4pt}
\small
\begin{tabularx}{\textwidth}{X l l}
\toprule
\textbf{Meta} & \textbf{Alvo} & \textbf{Resultado} \\
\midrule
T\'{o}picos essenciais de C\'{a}lculo~1 da UnB & $\geq 4$ & Atingido \\
Quest\~{o}es com explica\c{c}\~{a}o anal\'{i}tica completa & $\geq 20$ & 24 (superado) \\
Testes unit\'{a}rios automatizados & Cobertura de servi\c{c}os & 17 testes (100\% sucesso) \\
Build de produ\c{c}\~{a}o & Compila\c{c}\~{a}o sem erros & Atingido \\
Pacote mobile via Capacitor & Android/iOS & Sincronizado \\
Documenta\c{c}\~{a}o t\'{e}cnica & Todos os artefatos & 12 documentos \\
\bottomrule
\end{tabularx}
\end{table}

% ================================================================
\section{Documenta\c{c}\~{a}o Produzida}
\label{sec:documentacao}

\begin{table}[H]
\centering
\caption{Artefatos de documenta\c{c}\~{a}o do projeto}
\label{tab:documentacao}
\vspace{4pt}
\small
\begin{tabularx}{\textwidth}{l X}
\toprule
\textbf{Documento} & \textbf{Conte\'{u}do} \\
\midrule
Vis\~{a}o do Produto & Declara\c{c}\~{a}o de problema, posicionamento e objetivos estrat\'{e}gicos \\
Canvas MVP & S\'{i}ntese do produto m\'{i}nimo vi\'{a}vel em p\'{a}gina \'{u}nica \\
Personas & 3 perfis detalhados (Calouro, Veterana, Monitor) \\
Requisitos (RF/RNF) & 15 requisitos funcionais e 8 n\~{a}o-funcionais \\
Backlog & 18 hist\'{o}rias de usu\'{a}rio em 6 \'{e}picos \\
Arquitetura & Camadas, diagrama de fluxo e padr\~{o}es de projeto \\
Diagramas UML & Casos de uso, classes, componentes e sequ\^{e}ncias \\
Metodologia & Scrum adaptado e pr\'{a}ticas de engenharia \\
Cronograma & 5 sprints com marcos e entreg\'{a}veis \\
Qualidade & Ferramental, crit\'{e}rios de aceite e m\'{e}tricas \\
Riscos & Mapeamento e plano de mitiga\c{c}\~{a}o \\
Guia de Contribui\c{c}\~{a}o & Padr\~{o}es de c\'{o}digo e fluxo de trabalho \\
\bottomrule
\end{tabularx}
\end{table}

% ================================================================
\section{Conclus\~{a}o}
\label{sec:conclusao}

O desenvolvimento do BeeGe ao longo de 30 dias representou n\~{a}o apenas um desafio t\'{e}cnico, mas sobretudo um \textbf{salto pessoal} na dire\c{c}\~{a}o do empreendedorismo. O projeto demonstrou que \'{e} poss\'{i}vel, com planejamento, metodologia \'{a}gil \cite{schwaber2020} e as ferramentas certas, construir um MVP funcional e robusto em um prazo curto.

O resultado \'{e} uma plataforma que:

\begin{itemize}[leftmargin=2cm]
  \item Resolve um problema real enfrentado por milhares de estudantes universit\'{a}rios.
  \item Foi constru\'{i}da com arquitetura extens\'{i}vel e pronta para escalar como SaaS.
  \item Possui documenta\c{c}\~{a}o completa de engenharia de software \cite{pressman2021}.
  \item Foi validada com testes automatizados e builds de produ\c{c}\~{a}o.
  \item Conta com identidade visual pr\'{o}pria e mascote pedag\'{o}gica.
\end{itemize}

Mais do que um projeto t\'{e}cnico, o BeeGe \'{e} a materializa\c{c}\~{a}o da coragem de empreender --- um passo que nasceu dentro do programa NaPr\'{a}tica e que agora est\'{a} pronto para se tornar uma empresa real, capaz de transformar a maneira como estudantes de todo o Brasil revisam e aprendem C\'{a}lculo~1.

\vspace{1cm}

\begin{center}
\rule{0.5\textwidth}{0.4pt}\\[0.5cm]
\textit{``Me desafiei a fazer algo que eu sempre quis, mas n\~{a}o tinha coragem: empreender.''}\\[0.3cm]
\textbf{--- Gabriel Monteiro, 2026}\\[0.5cm]
\rule{0.5\textwidth}{0.4pt}
\end{center}

% ================================================================
\addcontentsline{toc}{section}{Refer\^{e}ncias}
\begin{thebibliography}{99}

\bibitem{stewart2016}
STEWART, J. \textit{C\'{a}lculo}. 8. ed. S\~{a}o Paulo: Cengage Learning, 2016.

\bibitem{ries2011}
RIES, E. \textit{The Lean Startup}. New York: Crown Business, 2011.

\bibitem{pressman2021}
PRESSMAN, R. S.; MAXIM, B. R. \textit{Engenharia de Software: uma abordagem profissional}. 9. ed. Porto Alegre: AMGH, 2021.

\bibitem{sommerville2019}
SOMMERVILLE, I. \textit{Engenharia de Software}. 10. ed. S\~{a}o Paulo: Pearson, 2019.

\bibitem{schwaber2020}
SCHWABER, K.; SUTHERLAND, J. \textit{The Scrum Guide}. 2020. Dispon\'{i}vel em: \url{https://scrumguides.org/scrum-guide.html}.

\bibitem{supabase2024}
SUPABASE INC. \textit{Supabase Documentation}. 2024. Dispon\'{i}vel em: \url{https://supabase.com/docs}. Acesso em: out. 2026.

\bibitem{ionic2024}
IONIC TEAM. \textit{Ionic Framework Documentation}. 2024. Dispon\'{i}vel em: \url{https://ionicframework.com/docs}. Acesso em: out. 2026.

\bibitem{katex2024}
KHAN ACADEMY. \textit{KaTeX --- The fastest math typesetting library for the web}. 2024. Dispon\'{i}vel em: \url{https://katex.org/}. Acesso em: out. 2026.

\bibitem{capacitor2024}
IONIC TEAM. \textit{Capacitor: Cross-platform Native Runtime for Web Apps}. 2024. Dispon\'{i}vel em: \url{https://capacitorjs.com/docs}. Acesso em: out. 2026.

\bibitem{basili1994gqm}
BASILI, V. R.; CALDIERA, G.; ROMBACH, H. D. The Goal Question Metric Approach. In: \textit{Encyclopedia of Software Engineering}. John Wiley \& Sons, 1994.

\end{thebibliography}

\newpage
\appendix
\section{Declara\c{c}\~{a}o de Uso de Intelig\^{e}ncia Artificial}
Em conformidade com as boas pr\'{a}ticas de transpar\^{e}ncia acad\^{e}mica, o autor declara o uso de ferramentas de IA generativa nas seguintes atividades:
\begin{itemize}
    \item \textbf{Ferramentas Utilizadas:} Antigravity / Gemini 3.8 (Google DeepMind).
    \item \textbf{Consolida\c{c}\~{a}o por Tarefas:} aux\'{i}lio na gera\c{c}\~{a}o de c\'{o}digo boilerplate para componentes Ionic/React; suporte na formata\c{c}\~{a}o de documenta\c{c}\~{a}o t\'{e}cnica em Markdown e \LaTeX{}; e apoio na estrutura\c{c}\~{a}o de testes automatizados com Vitest.
    \item \textbf{Confer\^{e}ncia Cr\'{i}tica pelo Autor:} todo c\'{o}digo-fonte, dados e conte\'{u}dos gerados foram integralmente revisados, testados e validados pelo autor, garantindo a fidelidade t\'{e}cnica e a corretude do produto final.
\end{itemize}

\end{document}
